\subsection{A new variant for Stembridge's alternating tableaux}
\label{sec:Stembridge}

In this section we present a variation of Sundaram's bijection for
alternating tableaux and permutations.

Recall that a staircase is a vector with weakly decreasing integer
entries. The \Dfn{positive part} of the staircase is the partition
obtained by removing all entries less than or equal to zero. The
\Dfn{negative part} of the staircase is the partition obtained by
removing all entries greater than or equal to zero, removing the
signs of the remaining entries and reversing the sequence.
\begin{defi}\label{defn:P}
 Let $\alt=(\wgt^0,\wgt^1,\ldots,\wgt^{2r})$ be an alternating
 tableau. The associated growth diagram $\Grm(\alt)$ is an
 $r\times r$ square of cells, obtained as follows:
 \begin{itemize}\setlength{\itemsep}{0pt}
 \item[P1] Label the north east corners of the cells on the main diagonal and
 the first superdiagonal from the top-left to the bottom-right
 with the staircases in $\alt$.
 \item[P2] Apply the backward rules on the positive parts of the
 staircases to determine which cells below the diagonal contain a
 cross.
 \item[P3] Use the backward rules (rotated by $180\degree$) on the
 negative parts of the staircases to determine which cells above
 the diagonal contain a cross.
 \end{itemize}

 Let $\Perm(\alt)$ be the partial permutation mapping $i$ to $j$ for
 every cross in column $i$ and row $j$ of $\Grm(\alt)$, and let
 $\big(\Perm_P(\alt), \Perm_Q(\alt)\big)$ be the pair of partial
 standard Young tableaux corresponding to the sequence of partitions
 along the bottom and the right border of $\Grm(\alt)$, respectively.
\end{defi}
